\documentclass{article}
\usepackage[utf8]{inputenc}
\usepackage{hyperref}
\usepackage{listings}
\usepackage{xcolor}

\definecolor{codegreen}{rgb}{0,0.6,0}
\definecolor{codegray}{rgb}{0.5,0.5,0.5}
\definecolor{codepurple}{rgb}{0.58,0,0.82}
\definecolor{backcolour}{rgb}{0.95,0.95,0.92}

\lstdefinestyle{mystyle}{
    backgroundcolor=\color{backcolour},   
    commentstyle=\color{codegreen},
    keywordstyle=\color{magenta},
    numberstyle=\tiny\color{codegray},
    stringstyle=\color{codepurple},
    basicstyle=\ttfamily\footnotesize,
    breakatwhitespace=false,         
    breaklines=true,                 
    captionpos=b,                    
    keepspaces=true,                 
    numbers=left,                    
    numbersep=5pt,                  
    showspaces=false,                
    showstringspaces=false,
    showtabs=false,                  
    tabsize=2
}

\lstset{style=mystyle}

\title{Frontend Implementation: Intelligent Restaurant Management System}
\author{Software Architecture Documentation}
\date{April 2026}

\begin{document}

\maketitle

\section{Frontend Structure}
The project follows a \textbf{Feature-Driven Architecture}, where each business domain is encapsulated within its own module. This promotes high cohesion and low coupling.

\subsection{Directory Layout}
\begin{lstlisting}[language=bash]
intelligent-restaurant-fe/
├── app/                  # Next.js App Router (Pages & Layouts)
│   ├── (admin)/          # Admin-facing pages (Analytics)
│   ├── (auth)/           # Authentication pages (Login)
│   ├── (chef)/           # Chef-facing pages (KDS)
│   ├── (customer)/       # Customer-facing pages (Menu)
│   └── (staff)/          # Staff-facing pages (Tables)
├── components/           # Shared UI components (shadcn/ui)
├── features/             # Business Logic (Feature-Driven)
│   ├── [feature-name]/
│   │   ├── components/   # Reusable feature-specific UI
│   │   ├── config/       # Domain types & Zod schemas
│   │   └── data-access/  # API interfaces, classes, and Query hooks
├── lib/                  # Shared utilities and global config
└── providers/            # Global context providers (Realtime, Query)
\end{lstlisting}

\subsection{Key Module Components}
\begin{itemize}
    \item \textbf{Providers}: Located at the root level to avoid confusion with route paths. These provide global functionality like TanStack Query and Realtime event handling.
    \item \textbf{Config}: Centralizes domain definitions using TypeScript and Zod. This acts as the "Single Source of Truth" for data structures.
    \item \textbf{Data-Access}:
    \begin{itemize}
        \item \textbf{Interfaces}: Define the contract for API operations (Dependency Inversion).
        \item \textbf{Concrete Classes}: Implementations for \texttt{Mock} (localStorage) and \texttt{Real} (REST API) environments.
        \item \textbf{Queries}: Encapsulate TanStack Query hooks to manage caching and server-state.
    \end{itemize}
    \item \textbf{Components}: Functional UI parts that consume data through feature-specific hooks.
\end{itemize}

\section{SOLID Principles Application}

\subsection{S - Single Responsibility Principle (SRP)}
Each module and file has a dedicated responsibility. For instance, \texttt{*.api.ts} handles only the network layer, while \texttt{*.queries.ts} manages server state and caching.

\subsection{O - Open/Closed Principle (OCP)}
The system is open for extension but closed for modification through \textbf{Interface-Driven API}. New environments can be added by creating new classes that implement the existing API interface without touching UI logic.

\subsection{L - Liskov Substitution Principle (LSP)}
Our Mock/Real API pattern is a direct implementation of LSP. Both \texttt{MockMenuApi} and \texttt{RealMenuApi} strictly implement \texttt{IMenuApi}, allowing them to be substituted at runtime without affecting the application's correctness.

\subsection{I - Interface Segregation Principle (ISP)}
We use \textbf{Feature-Segregated Interfaces} (\texttt{IAuthApi}, \texttt{IMenuApi}, etc.). Features only depend on the interfaces they actually use, reducing coupling and cognitive load.

\subsection{D - Dependency Inversion Principle (DIP)}
High-level modules depend on \textbf{Abstractions} (Interfaces), not low-level details. Pages and components use hooks that depend on the API Interface, while the specific implementation is injected at the API level.

\end{document}
